\section {Barrier using semaphores (Programming)}

Implement a barrier for $n$ threads using semaphores.
A thread that calls \verb|barrier()| must sleep until all $n$ threads have called it, and then all of them continue.
The barrier must be reusable: the same threads will call it again and again.

\subsection{Solution}

We use three semaphores: \verb|mutex| protects the counter, \verb|sleep| is where the threads wait and
\verb|done| lets the last thread know that everybody woke up.
\verb|up(s, k)| and \verb|down(s, k)| add and subtract $k$ from a semaphore (as \verb|semop(2)| does).

\begin{verbatim}
init(n):
    set(mutex, 1)
    set(sleep, 0)
    set(done, 0)
    count = 0

barrier():
    down(mutex, 1)
    count += 1
    if count == n:
        # the last thread to arrive wakes everybody up
        count = 0
        up(sleep, n-1)
        down(done, n-1)
        up(mutex, 1)
    else:
        # order matters: release the mutex and only then sleep
        up(mutex, 1)
        down(sleep, 1)
        up(done, 1)
\end{verbatim}

A sleeping thread must release the mutex before it goes to sleep, otherwise no other thread could ever reach the barrier.

Why does the last thread wait on \verb|done| while holding the mutex?
Without it, a fast thread that was just woken could run ahead, call \verb|barrier()| again for the next round and grab one of the
$n-1$ wake up tokens of the current round. It would then pass the barrier without waiting for the next round, and one of the
threads of the current round would stay asleep. Holding the mutex until all $n-1$ threads have taken their tokens means that no
thread can start the next round before the current round is over.
